% Generated draft: requires mathematical and source review.
\documentclass[11pt]{article}
\usepackage[margin=25mm]{geometry}
\usepackage{fontspec}
\setmainfont{Latin Modern Roman}
\usepackage{amsmath,amssymb,mathtools,unicode-math}
\setmathfont{Latin Modern Math}
\usepackage{enumitem,needspace,xurl,hyperref}
\setlength{\parindent}{0pt}
\setlength{\parskip}{5pt}
\setlength{\emergencystretch}{3em}
\newcommand{\N}{\mathbb{N}}
\newcommand{\Z}{\mathbb{Z}}
\newcommand{\Q}{\mathbb{Q}}
\newcommand{\R}{\mathbb{R}}
\newcommand{\C}{\mathbb{C}}
\newcommand{\F}{\mathbb{F}}
\newcommand{\A}{\mathbb{A}}
\newcommand{\PP}{\mathbb{P}}
\newcommand{\E}{\mathbb{E}}
\newcommand{\eps}{\varepsilon}
\newcommand{\dd}{\,\mathrm{d}}
\begin{document}
\section*{UnsolvedMath ID: 30002171}
\textbf{Source problem code:} catalogue entry
\subsection{Definitions and mathematical statement}
Let \(\widetilde Z\) denote the spectral-curve partition function referred to in the supplied record, and write \(\mathbf t=(t_i)_{i\in I}\) for its intended deformation parameters, where \(I\) is infinite. The record does not specify the spectral-curve data, the construction or normalization of \(\widetilde Z\), the parameter space, the coefficient ring for formal expansions, or the intended Hirota hierarchy.

\begin{enumerate}
\item Does \(\widetilde Z(\mathbf t)\) formally satisfy the Hirota equations in the infinitely many deformation parameters \(\mathbf t\)? The hierarchy, formal expansion convention, and class of admissible spectral curves must be specified before this question determines a concrete system of identities or a universally quantified conjecture.
\item For a knot \(K\), let \(C_K\) denote the intended curve defined by its \(A\)-polynomial, and let \(\iota\) denote the involution referred to in the source. What is the knot-theoretic meaning of the observed genus-zero condition for the quotient
\[
D_K=C_K/\langle\iota\rangle?
\]
If the intended genus is geometric genus, this condition would be expressed as
\[
g\bigl(\overline{D_K}^{\,\mathrm{norm}}\bigr)=0,
\]
where \(\overline{D_K}^{\,\mathrm{norm}}\) denotes the normalization of a projective completion of the intended quotient curve. This is a conditional interpretation: the source does not fix the ambient field, curve components, quotient construction, genus convention, involution, or class of knots for which the observation is intended. In particular, it does not supply a claim that this condition holds for every knot.
\end{enumerate}

The supplied background describes the first question as a nonperturbative topological-recursion conjecture and also mentions a vanishing invariant in connection with the second question. Neither that construction nor that invariant is defined in the record, so no identification of the invariant with the displayed genus is asserted.
\subsection{Short English statement}
The record asks whether an unspecified spectral-curve partition function satisfies a Hirota hierarchy in infinitely many deformation parameters, and what the observed genus-zero quotient of a knot's \(A\)-polynomial curve signifies topologically. The constructions, hierarchy, involution, and scope of the observation require clarification before a fully precise formulation is possible.
\subsection{Sources}
\begin{itemize}
\item Supplied local UnsolvedMath registry record, permanent ID \(30002171\), exact source problem code \texttt{catalogue entry}, titled \emph{Hirota Equations and Knot Meaning of Spectral Curves}, in the Oberwolfach Reports Open Problems collection. No individual proposer is supplied; the proposed year is recorded as 2012.
\item Supplied report citation: \emph{Low-Dimensional Topology and Number Theory} (2013), Oberwolfach Reports, pp.\ 2541--2596, DOI \texttt{10.4171/owr/2012/42}. Source URL and supplied literature-source link, labeled ``OWR report 42'': \url{https://doi.org/10.4171/owr/2012/42}. The report is supplied as provenance for workshop open-problem discussions; no precise problem page or passage is provided.
\item Literature assessment dated \texttt{2026-08-21}: the supplied curation labels the problem open, describes a nonperturbative topological-recursion conjecture, and states that no general knot-theoretic interpretation of the vanishing invariant has since been established. These are supplied dated curation claims, not independently verified literature or mathematical conclusions. The background expressly requests independent verification of time-sensitive status and cited claims.
\item The registry labels the statement \texttt{corrected\_verified} and reports a manual semantic, typographic, and LaTeX review of its cleanup. This is supplied editorial metadata, not independent mathematical verification of the formulation or conjecture.
\end{itemize}
\end{document}
